\documentclass[11pt,a4paper]{article}
\usepackage[margin=1in]{geometry}
\usepackage{amsmath,amssymb}
\usepackage{booktabs}
\usepackage{graphicx}
\usepackage{hyperref}
\usepackage{xcolor}
\title{Path to 1\,\textcent/kWh Fusion Energy}
\author{1cFE Project}
\date{March 2026}

\newcommand{\target}{\$10/\text{MWh}}

\begin{document}
\maketitle

\begin{abstract}
We decompose the \$10/MWh ($1\,\textcent$/kWh) levelized cost of
electricity target for fusion power into its constituent requirements,
using the \textsc{1costingfe} techno-economic model.  Starting from a
baseline of \$43/MWh for a 1\,GWe NOAK pB11 mirror plant, we identify
the lever stack required to close the $4.3\times$ gap.  We examine two
bounding cases: (i)~a free fusion core, which isolates the irreducible
balance-of-plant floor, and (ii)~a fully costed core, which represents
the engineering challenge.  In both cases, direct energy conversion
(DEC) emerges as a transformative lever by eliminating the steam cycle
that dominates non-core costs.  No single lever reaches the target;
the path requires simultaneous advances in plant scale, availability,
construction speed, financing, and---for charged-particle-dominated
fuels---DEC.
\end{abstract}

\tableofcontents


%% ====================================================================
\section{Baseline: Where Are We Today?}

The cheapest configuration in the current model is a 1\,GWe NOAK pB11
mirror plant at standard financial assumptions (7\% WACC, 85\%
availability, 30-year life, 6-year construction):

\begin{center}
\begin{tabular}{lr}
\toprule
Metric & Value \\
\midrule
LCOE & \$42.9/MWh \\
Overnight cost & \$3{,}548/kW \\
CAS22 (reactor equipment) & \$1{,}563M \\
CAS21 (buildings) & \$332M \\
CAS23--26 (BOP) & \$421M \\
CAS70 (annual O\&M) & \$24M/yr \\
\bottomrule
\end{tabular}
\end{center}

Capital charges (CAS90) account for $\sim$90\% of LCOE, with O\&M
contributing $\sim$10\% and fuel $<$1\%.  The gap to the \target{}
target is \$33/MWh---a $4.3\times$ reduction.  Because capital
dominates, the path to \target{} is overwhelmingly a capital cost
problem.

CAS21 buildings are priced at industrial/enhanced-industrial grade
appropriate for pB11 under 10~CFR~Part~30---not nuclear-grade
(Part~50).  pB11 buildings at \$308/kW reflect zero hot cell, zero
tritium infrastructure, standard HVAC, and industrial construction
throughout.


%% ====================================================================
\section{Free Fusion Core: The Balance-of-Plant Floor}

To isolate the irreducible non-core costs, we zero out all reactor
plant equipment (CAS22 = 0, CAS27 = 0).  Even with a completely free
heat source, the LCOE floor is \textbf{\$18/MWh (1.8\,\textcent/kWh)}
for pB11.  The fusion core accounts for 58\% of baseline LCOE.

\subsection{Closing the Gap}

\begin{center}
\begin{tabular}{lrr}
\toprule
Cumulative scenario & LCOE (\$/MWh) & O/N (\$/kW) \\
\midrule
Free core (baseline) & 18 & 1{,}231 \\
+ 2\,GWe, 95\%, 3\% WACC, 3yr, 50yr & \textbf{7.3} & 854 \\
\bottomrule
\end{tabular}
\end{center}

\subsection{Effect of Direct Energy Conversion}

\begin{center}
\begin{tabular}{lrr}
\toprule
Scenario & LCOE (\$/MWh) & O/N (\$/kW) \\
\midrule
Steam: 2\,GWe, 95\%, 3\% WACC, 3yr, 50yr & 7.3 & 854 \\
DEC: 2\,GWe, 95\%, 3\% WACC, 3yr, 50yr & \textbf{5.7} & 524 \\
DEC: 5\,GWe, 95\%, 2\% WACC, 3yr, 50yr & \textbf{3.9} & 448 \\
\bottomrule
\end{tabular}
\end{center}

At mega-scale with sovereign financing, the DEC floor reaches
\$3.9/MWh---dominated almost entirely by O\&M staff costs.


%% ====================================================================
\section{Fully Costed Fusion Core}

With a realistic fusion core (CAS22 computed from first principles),
the baseline LCOE is \$43/MWh for pB11 mirror.

\subsection{Combined Scenario Stacking}

\begin{center}
\begin{tabular}{lrr}
\toprule
Cumulative scenario & LCOE (\$/MWh) & O/N (\$/kW) \\
\midrule
Baseline (1\,GWe, 7\% WACC, 30yr) & 42.9 & 3{,}548 \\
+ 2\,GWe, 95\%, 3yr, 3\% WACC, 50yr & \textbf{12.5} & 1{,}982 \\
\bottomrule
\end{tabular}
\end{center}

With all levers stacked and a steam cycle, the fully costed pB11
mirror reaches \$12.5/MWh---close but still above target.

\subsection{Effect of Direct Energy Conversion}

\begin{center}
\begin{tabular}{lrrrl}
\toprule
Scenario & $f_{\text{dec}}$ & LCOE (\$/MWh) & O/N (\$/kW) & Notes \\
\midrule
Baseline (steam only) & 0\% & 42.9 & 3{,}548 & \\
DEC at baseline & 90\% & 35.8 & 2{,}896 & BOP reduced \\
DEC + 2\,GWe, 95\%, 3yr, 3\%, 50yr & 90\% & 10.4 & 1{,}526 & At target \\
DEC + 5\,GWe, 95\%, 3yr, 2\%, 50yr & 90\% & 6.1 & 1{,}026 & Below target \\
\bottomrule
\end{tabular}
\end{center}

\textbf{Key finding:} With DEC, the fully costed pB11 mirror reaches
\$10.4/MWh at moderate scale (2\,GWe) and favorable-but-not-extreme
financing (3\% WACC, 50-yr life, 3-yr build)---essentially at target.
At mega-scale (5\,GWe, 2\% WACC), \$6.1/MWh.

\subsection{Automation}

The remaining pB11-specific lever is \textbf{automated operations}---
reducing the 59-person baseline staff to $\sim$30~FTE via a
data-center operating model (remote monitoring, AI predictive
maintenance, robot-accessible plant with zero activation barriers).

\begin{center}
\begin{tabular}{lrr}
\toprule
Scenario & LCOE (\$/MWh) & O/N (\$/kW) \\
\midrule
DEC + moderate (59~FTE) & 10.4 & 1{,}526 \\
DEC + moderate + automated (30~FTE) & \textbf{8.8} & 1{,}526 \\
\bottomrule
\end{tabular}
\end{center}

Automation buys \$1.6/MWh---the difference between ``at target'' and
``below target with margin.''


%% ====================================================================
\section{Summary of Paths}

\begin{center}
\begin{tabular}{clccp{5.5cm}}
\toprule
Path & Description & LCOE & Core & Key requirement \\
\midrule
1 & Steam, free core, all levers & \$7.3 & Free &
  2\,GWe, 95\%, 3\% WACC, 3yr, 50yr \\
2 & DEC, free core, moderate & \$5.7 & Free &
  2\,GWe, 95\%, 3\% WACC, 3yr, 50yr \\
3 & DEC, free core, mega-scale & \$3.9 & Free &
  5\,GWe, 95\%, 2\% WACC, 3yr, 50yr \\
\midrule
4 & Steam, full core, all levers & \$12.5 & Full &
  2\,GWe, 95\%, 3\% WACC, 3yr, 50yr \\
5 & DEC, full core, moderate & \$10.4 & Full &
  2\,GWe, 95\%, 3\% WACC, 3yr, 50yr \\
6 & DEC, full core, mega-scale & \$6.1 & Full &
  5\,GWe, 95\%, 2\% WACC, 3yr, 50yr \\
\textbf{7} & \textbf{DEC + automated, full core} & \textbf{\$8.8} & \textbf{Full} &
  \textbf{2\,GWe, 95\%, 3\% WACC, 3yr, 50yr, 30\,FTE} \\
\bottomrule
\end{tabular}
\end{center}

\subsection{Required Conditions}

All paths to \target{} require:

\begin{enumerate}
\item \textbf{Aneutronic fuel (pB11).}  DT adds $\sim$\$1.7B/GWe in
  neutron infrastructure.  Even with a free DT core, the floor is
  $\sim$\$40/MWh vs \$18/MWh for pB11.
\item \textbf{Large plant ($\geq$2\,GWe net).}  Fixed costs must be
  spread across many MWh.
\item \textbf{High availability ($\geq$95\%).}
\item \textbf{Low cost of capital ($\leq$3\% real WACC).}
\item \textbf{Fast construction ($\leq$3\,yr).}
\end{enumerate}

\subsection{The DEC Dividend}

DEC is the discriminating lever between ``close but not quite'' and
``at target.''  For pB11:
\begin{itemize}
\item Without DEC: \$12.5/MWh at moderate scale (25\% above target).
\item With DEC: \$10.4/MWh at moderate scale (at target).
\item With DEC + automation: \$8.8/MWh (below target with margin).
\end{itemize}


%% ====================================================================
\section{What Definitely Will Not Work}

\begin{itemize}
\item \textbf{7\% WACC:} Requires $<$\$500/kW overnight---below gas
  turbines, physically implausible.
\item \textbf{1\,GWe scale:} Fixed costs too high even at 2\% WACC.
\item \textbf{DT fuel:} Free-core floor ($\sim$\$40/MWh) is close to
  pB11's fully costed result (\$43/MWh).
\item \textbf{Any single lever alone:} Deepest reach is $\sim$\$22/MWh
  (5\,GWe scale).
\end{itemize}


%% ====================================================================
\section{Conclusions}

\begin{enumerate}
\item \textbf{The fusion core is necessary but not sufficient.}  Even
  a free core leaves an \$18/MWh floor.

\item \textbf{DEC is the highest-value technology lever for pB11.}
  It provides $\sim$\$2/MWh of headroom at moderate scale.

\item \textbf{Scale matters more than unit cost reduction.}  Going
  from 1 to 2\,GWe buys more LCOE reduction than any other single
  lever.

\item \textbf{Financial conditions are as important as engineering.}
  Government-backed financing at 3\% WACC is worth $\sim$\$15/MWh
  versus commercial 7\% WACC.

\item \textbf{The path is pB11 + DEC + scale + favorable finance.}
  Path~5 produces \$10.4/MWh with a fully costed reactor.  Path~7
  (+ automation) reaches \$8.8/MWh.
\end{enumerate}

The key physics challenge is not reaching ignition---it is reaching
ignition in a configuration that supports DEC, high availability,
and multi-GWe deployment at low unit cost.

\end{document}
